\subsection{The entire beta-trivial leading term in degree eleven}
\label{subsec:eleven-completeness}

Willwacher's rational dimension bound identifies the complete leading
term removed by Theorem~\ref{thm:eleven-mixed-mixed-reduction}. The
arithmetic unit then gives its exact integral lattice.

For $g\in\ker\lambda$ whose dyadic Magnus word has no term below
degree $n$, write $\operatorname{in}_n(g)$ for its degree-$n$
homogeneous term. Write $F^n$ for this closed word filtration.

\begin{lemma}\label{lem:full-gt-leading-term}
If $g$ is a full-GT element in $F^n$, with $n\ge3$, then
\[
 \operatorname{in}_n(g)\in
 \mathfrak{grt}_1(\mathbf Q_2)_n
       \cap\mathbf Z_2\langle S,T\rangle_n.
\]
Furthermore
\[
 \mathfrak{grt}_1(\mathbf Q_2)_n
  =\mathfrak{grt}_1(\mathbf Q)_n\otimes_{\mathbf Q}\mathbf Q_2.
\]
\end{lemma}

\begin{proof}
Project the profinite word and its full pentagon relation to the
pro-$2$ completions, then to the rational prounipotent completions.
The pro-$\ell$ to prounipotent passage is explicitly recalled in
Furusho's proof of the pro-$\ell$ corollary
\cite[\S3, p.~554, immediately before Corollary~12]{Furusho2010}.
Set $x=e^S$, $y=e^T$ in the free prounipotent group. The leading
term is unchanged by this substitution, because
$e^S-1=S+O(S^2)$ and similarly for $T$. Since the image of $f_g$
is group-like, its first nonzero term is a primitive element of
the free associative Hopf algebra in characteristic zero, hence
a homogeneous free Lie polynomial.

Taking the first nonzero degree in the full pentagon gives
\begin{align*}
 q(t_{12},t_{23}+t_{24})+q(t_{13}+t_{23},t_{34})
 ={}&q(t_{23},t_{34})\\
   &+q(t_{12}+t_{13},t_{24}+t_{34})+q(t_{12},t_{23}),
\end{align*}
where $q=\operatorname{in}_n(g)$ and the equality is in the
Drinfeld--Kohno Lie algebra $\mathfrak t_4(\mathbf Q_2)$.
Indeed, the degree-one class of a pure braid generator is
$t_{ij}$, a product of generators has the sum of these classes,
and a product of two nonconstant $f_g$ terms has degree at least
$2n>n$. This is the associated graded form of the usual
prounipotent pure braid realization. Equivalently it follows by
the associator transport used in the proof of
\cite[proof of Theorem~2, p.~554]{Furusho2010}: that transport preserves
linear terms, so it preserves the displayed leading equation.

Duality and the multiplier-one hexagon similarly give the
antisymmetry and three-cycle relations. In characteristic zero
the cabling relation is redundant, as recalled with Drinfeld's
Proposition~5.7 in \cite[p.~1, after equation~(6)]{Willwacher}.
Thus $q$ satisfies precisely the defining equations (1)--(4)
of that source. One may also use the pentagon implication
\cite[Theorem~3, pp.~547--548]{Furusho2010}, together with the
same cabling redundancy. The hypotheses on degree exclude all
linear and quadratic exceptional terms. This proves the first
assertion. It is a necessary leading-term assertion only; no
surjectivity from profinite GT to a prounipotent group is used.

In fixed degree the free Lie and infinitesimal braid Lie algebras
have finite bases over $\mathbf Q$, and all defining equations
have rational, indeed integral, coefficients. This integral
matrix description is noted explicitly in
\cite[\S1.1, p.~2]{Willwacher}. Choose these bases and write the
finite defining map as a matrix $A_n$ over $\mathbf Q$.
Its rank is unchanged after extension to $\mathbf Q_2$, since a
minor is nonzero over one characteristic-zero field if and only
if it is nonzero over the other. Equivalently, tensoring the
kernel sequence with the field extension is exact. Hence
$\ker(A_n\otimes\mathbf Q_2)=(\ker A_n)\otimes\mathbf Q_2$,
which is the asserted base-change identity.
\end{proof}

\begin{proposition}\label{prop:eleven-complete-leading-line}
Let $L_8$ and $P_3$ have the normalizations in
Theorem~\ref{thm:eleven-mixed-mixed-reduction}, and put
\[
 V_{11}=\frac1{11}\{P_3,L_8\}.
\]
Then $V_{11}$ is an integral $2$-adic Lie polynomial with
$[S^6TST^3]V_{11}=1$, and
\[
 \operatorname{in}_{11}(g)=b_{11}(g)V_{11}\quad(g\in E).
\]
In particular
\[
 \operatorname{in}_{11}(E)
 =\operatorname{in}_{11}(E\cap H)=\mathbf Z_2V_{11},
 \qquad E'=E\cap F^{12}=C_\infty\cap F^{12}.
\]
Thus every dyadic word belonging to $E'$ starts in degree twelve.
All arithmetic factorizations, relative quotient identifications,
and character-kernel conclusions of the strict reduction remain
valid for this deeper subgroup.
\end{proposition}

\begin{proof}
Willwacher's dimension bound
\cite[Theorem~1.1, p.~2]{Willwacher} gives
\[
 \dim_{\mathbf Q_2}\mathfrak{grt}_1(\mathbf Q_2)_{11}
 =\dim_{\mathbf Q}\mathfrak{grt}_1(\mathbf Q)_{11}\le d_{11}=2.
\]
Here the equality of dimensions is
Lemma~\ref{lem:full-gt-leading-term}. To check the last integer
directly, the free Lie algebra on generators of odd weights
$3,5,7,\ldots$ has only two possibilities in weight eleven:
the generator of weight eleven, and the multidegree with two
weight-three letters and one weight-five letter. The latter
has one-dimensional span, represented by
$[\sigma_3,[\sigma_3,\sigma_5]]$; brackets of two generators
have even weight and brackets of at least four have weight at
least twelve. We need no injectivity theorem for free generators.

We supply the two independent directions using arithmetic
elements. The preceding strict reduction already supplies
$w_{11}\in E\cap H$ with
\[
 [S^{10}T]\operatorname{in}_{11}(w_{11})=0,\qquad
 b_{11}(w_{11})=99a^2u\in\mathbf Z_2^\times.
\]
The first equality follows because its beta multiplier is one.

For the second direction, let
\[
 v_5=\Delta_9\Delta_7\Delta_5\Delta_3\Delta_1(\sigma),
 \qquad \Delta_j(z)=hzh^{-1}z^{-t^j},\quad t_2=3.
\]
The uniform odd normalization gives $v_5\in C_4\cap H$ and
\[
 m_{11}(v_5)=
 \frac{2\prod_{j=0}^{4}(3^{11}-3^{2j+1})}{10!}
                     \tau_{11}(\sigma)
       \in2^{11}\mathbf Z_2^\times.
\]
Indeed the numerator product has valuation $18$, the extra
factor two adds one, and $v_2(10!)=8$; the moment is odd.
Here $C_4$ kills the depth-one weights $3,5,7,9$ but need not
kill the mixed terms of weights eight and ten.

Use the established beta-trivial arithmetic units $w_8$ and
$w_{10}$, with $m_8(w_8)$ and $m_{10}(w_{10})$ odd. First multiply $v_5$ by
a suitable profinite power of $w_8$ to kill $m_8$. The result
is in $D_8\cap C_4\cap H$, hence in $D_9$, since its
weight-nine depth-one coordinate is still zero. Next multiply
by a suitable power of $w_{10}$ to kill $m_{10}$. This gives
$z_{11}\in D_{10}\cap H$. Both correcting words have beta
multiplier one, so all beta coefficients of $z_{11}$ equal
those of $v_5$. In particular
\[
 [S^{10}T]\operatorname{in}_{11}(z_{11})
       =m_{11}(z_{11})=m_{11}(v_5)\ne0.
\]
The established degree-ten normalization ensures that
$z_{11}$ has no term below degree eleven.

By Lemma~\ref{lem:full-gt-leading-term}, both leading polynomials
belong to $\mathfrak{grt}_1(\mathbf Q_2)_{11}$. They are linearly
independent, as detected by $[S^{10}T]$ and $b_{11}$. This proves
that the dimension is exactly two and that the kernel of the
depth-one functional $[S^{10}T]$ is the one-dimensional line
spanned by $\operatorname{in}_{11}(w_{11})$.

Every $g\in E$ lies in that kernel, since $B_g=1$. Its leading
term is therefore a scalar multiple of $V_{11}$. The exact
coefficient identity
$[S^6TST^3]\{P_3,L_8\}=11$ identifies that scalar with
$b_{11}(g)$. Division by eleven is harmless in $\mathbf Z_2$,
so $V_{11}$ itself is integral. The arithmetic unit from the
strict reduction realizes all scalars in $\mathbf Z_2$ and
establishes both leading-image equalities. This is the integral
saturation step; it does not follow from rational dimension
alone.

It follows that $b_{11}(g)=0$ on $E$ if and only if its entire
degree-eleven term vanishes. This proves $E'=E\cap F^{12}$.
Finally $C_\infty\cap F^{12}\subset D_{10}$, because all the
earlier mixed and depth-one coordinates vanish when the word
starts in degree twelve. The last displayed equality follows.
\end{proof}

